\documentclass[11pt,a4paper]{article}
\usepackage[utf8]{inputenc}
\usepackage[T1]{fontenc}
\usepackage[brazil]{babel}
\usepackage[a4paper,margin=2.3cm]{geometry}
\usepackage{enumitem}
\usepackage{hyperref}
\usepackage{xcolor}
\setlength{\parindent}{0pt}
\setlength{\parskip}{5pt}
\sloppy
\setlength{\emergencystretch}{3em}
\hypersetup{colorlinks=true,urlcolor=blue,linkcolor=black}
\newcommand{\cabecalho}[2]{\begin{center}{\small \textbf{Sprint 2 -- #1}}\\[1mm]\textit{Inventário de Ativos e Vulnerabilidades}\\Matheus Moreli Durão Silva -- 12621CBS217\\Cibersegurança -- Bacharelado -- Noturno\end{center}\hrule\vspace{4mm}}
\newcommand{\bloco}[1]{\par\vspace{2mm}\textbf{#1}\par}
\newcommand{\codigo}[1]{\texttt{#1}}

\begin{document}\cabecalho{S2\_14 -- Consulta de Vulnerabilidades}{}
\bloco{Consulta local} \codigo{consultar\_vulnerabilidades} percorre todos os ativos e suas listas. Ela pode procurar CVE exato, prioridade e status. O retorno é uma lista de pares \codigo{(ativo, vulnerabilidade)}; nenhum atributo é alterado durante a busca.
\bloco{Evidência externa} A consulta local mostra fonte e data salvas no momento do cadastro. Como exemplo real, foi consultado em 24/09/2026 o registro \codigo{CVE-2021-44228} no NVD (\url{https://nvd.nist.gov/vuln/detail/CVE-2021-44228}), referente ao Apache Log4j e com CVSS 10.0. A referência externa não é substituída por uma lembrança ou por texto sem origem.
\bloco{Exemplo} Informar \codigo{CVE-2021-44228} retorna somente registros cujo CVE coincida exatamente, sem alterar o status ou prioridade. Antes do cadastro, deve-se confirmar se o ativo realmente usa o software afetado. Se a lista final estiver vazia, a interface mostra ``Nenhuma vulnerabilidade encontrada''.
\bloco{Como testar} O teste cadastra uma vulnerabilidade de exemplo e em seguida filtra pelo mesmo CVE. A quantidade esperada é um. Também se pode consultar um CVE inexistente e observar a mensagem controlada.
\bloco{Checklist atendido} CVE exato, filtros sem mutação, fonte/data exibidas, ausência tratada e procedimento de consulta externa documentado.
\end{document}
